\input format.tex

\usepackage{graphicx}
\graphicspath{{cores/}}

\usepackage{environ}
\usepackage{colortbl,array,booktabs}
\usepackage{tabularx}
\newcommand{\mytoprule}{\specialrule{0.1em}{0em}{0em}}
\newcommand{\mybottomrule}{\specialrule{0.1em}{0em}{0em}} 

\usepackage[nodisplayskipstretch]{setspace}

\colorlet{TablaBordeSuperior}{topcolor}
\colorlet{TablaBordeInferior}{topcolor}
\colorlet{TablaCentroSuperior}{blue!1}
\colorlet{TablaCentroInferior}{blue!20}
\colorlet{FuenteCabeceraTabla}{white}

\newcolumntype{M}[1]{>{\centering\arraybackslash}m{#1}}

\tcbset{rtab/.style={
freelance,
frame code={
 \path[top color=topcolor,bottom color=topcolor]
   ([yshift=-#1*(\baselineskip+2pt)]interior.north west) --
   ([yshift=-#1*(\baselineskip+2pt)]interior.north east) {[sharp corners]--
    ([yshift=3pt]interior.north east) --
    ([yshift=3pt]interior.north west)} -- cycle;

  },
interior code={},
 }
}

\newcommand\fuentecabecera[1]{\textcolor{black}{\textbf{#1}}}

\NewEnviron{RCtable}[1]
  {%
    \begin{tcolorbox}[#1]\BODY\end{tcolorbox}%
  }

\usepackage{colortbl}

\begin{document}

\vspace*{14mm}
%% 各章节
\setlength{\arrayrulewidth}{1pt}
\fontsize{9.3pt}{11pt}\selectfont
\color{gray2}

\begin{tctabularx}{tabularx={m{15.4cm}<{\centering}}}
\\[-6pt]
\cellcolor{topcolor}\raisebox{2.614pt}{\color{black70}\fontsize{11pt}{11pt}\selectfont\bf 主要有益菌检测结果}
\end{tctabularx}

\sepline

\noindent
\fontsize{8.8pt}{11pt}\selectfont
\arrayrulecolor{gray2}
\begin{tabular}{p{9.5cm}p{5.5cm}}
\parbox[l]{\hsize}{
\vspace*{5mm}
{
\tikz\draw[topcolor,fill=topcolor](0,0) rectangle(0.1,0.25);
\fontsize{10pt}{11pt}\selectfont\bf {乳杆菌属（Lactobacillus）}
}} & \\
\hspace*{1.5mm}
\begin{minipage}[t]{.6\textwidth}{
\vspace*{2mm}
\noindent
\begin{spacing}{1.5}
\fontsize{8pt}{11pt}\selectfont {乳杆菌，又叫乳酸菌，大家一定不陌生。很早以前，它就被人用来发酵生产酸奶和奶酪等食物，让食物既美味又营养。可不能小瞧了乳杆菌，增强免疫力、抵抗病菌、延缓衰老、促进营养吸收等都有它的功劳。当肠道内乳杆菌含量较高时，产生乳酸，调节肠道微环境，阻止致病菌侵入肠道；调节免疫，抑制肠道炎症的发生；促进营养物质吸收，从而有利于肠道健康。当肠道内乳杆菌含量较低时，肠道环境差，易发生肠道炎症、肠道问题甚至可能影响全身。这时记得适量补充酸奶或乳杆菌制剂哦，也可以搭配服用能促进乳杆菌生长的益生元。}
\end{spacing}
\vspace{3pt}
}
\end{minipage}
 & 
\hspace*{8mm}
\begin{minipage}[t]{.3\textwidth}{
\vspace*{2mm}
\tikz\draw[topcolor,fill=topcolor](0,0) rectangle(0.1,0.25);
{\fontsize{9pt}{11pt}\selectfont\bf 健康状态提示 \vspace{3pt}\\}
{\includegraphics[scale=1.0]{laughwithgreen.pdf} \vspace*{10mm}} \\
\tikz\draw[topcolor,fill=topcolor](0,0) rectangle(0.1,0.25);
{\fontsize{9pt}{11pt}\selectfont{\bf 检测结果评价} ~~~~正常 \vspace{3pt}\\}
{\fontsize{8pt}{11pt}\selectfont ~~~~正常抑制肠道炎症及调节肠道环境}
}
\end{minipage} \\
\mytoprule
\parbox[l]{\hsize}{
\vspace*{5mm}
{
\tikz\draw[topcolor,fill=topcolor](0,0) rectangle(0.1,0.25);
\fontsize{10pt}{11pt}\selectfont\bf {双歧杆菌属（Bifidobacterium）}
}} & \\
\hspace*{1.5mm}
\begin{minipage}[t]{.6\textwidth}{
\vspace*{2mm}
\noindent
\begin{spacing}{1.5}
\fontsize{8pt}{11pt}\selectfont {双歧杆菌是乳杆菌的好朋友，作用不比乳杆菌小：它既是肠道内的“清道夫”，也是调节肠道生态平衡的重要细菌。肠道内双歧杆菌含量较高时，表明肠道较为健康：它不仅能帮助增强肠道屏障，调节肠道环境；还能调节肠道免疫功能，有效抑制有害菌的胡作非为。此外，双歧杆菌还能够促进营养物质吸收，保持健康的身体状态。当肠道内双歧杆菌含量较低时，肠道可能进入不健康状态：肠道环境差、免疫调节能力低。这时需要补充双歧杆菌制剂或者益生元，增加双歧杆菌含量，优化肠道菌群组成。}
\end{spacing}
\vspace{3pt}
}
\end{minipage}
 & 
\hspace*{8mm}
\begin{minipage}[t]{.3\textwidth}{
\vspace*{2mm}
\tikz\draw[topcolor,fill=topcolor](0,0) rectangle(0.1,0.25);
{\fontsize{9pt}{11pt}\selectfont\bf 健康状态提示 \vspace{3pt}\\}
{\includegraphics[scale=1.0]{crywithorange.pdf} \vspace*{10mm}} \\
\tikz\draw[topcolor,fill=topcolor](0,0) rectangle(0.1,0.25);
{\fontsize{9pt}{11pt}\selectfont{\bf 检测结果评价} ~~~~偏低 \vspace{3pt}\\}
{\fontsize{8pt}{11pt}\selectfont ~~~~不利于调节免疫及肠道环境}
}
\end{minipage} \\
\mytoprule
\end{tabular}

\vspace*{3mm}
\noindent
~~~~\tikz\draw[green2,fill=green2](0,0) circle[radius=4pt]; \bf 绿色表示健康 \tikz\draw[orange2,fill=orange2](0,0) circle[radius=4pt]; \bf 橙色表示需要关注 \tikz\draw[red2,fill=red2](0,0) circle[radius=4pt]; \bf 红色表示为有风险

\end{document}
